%% shortname: Second-Order Stochastic Dominance and Introduction to General Equilibrium
\documentclass[11pt]{article}
\usepackage{microstyle}

\begin{document}

\lessonheader{19}{Chapter 6 (concluded): second-order stochastic dominance; Chapter 15: introduction to general equilibrium and the Edgeworth box}{}


\section{Reminder: first-order stochastic dominance}

$F$ first-order stochastically dominates (FOSD) $G$ if
\[ \int u(x)\,dF(x)\ \ge\ \int u(x)\,dG(x)
   \quad\text{for all non-decreasing utility functions } u(x) \tag{1} \]
\[ \Longleftrightarrow\qquad F(x)\ \le\ G(x) \qquad \forall\, x\in[0,\bar x]. \tag{2} \]

FOSD is very powerful, but $F(x)\le G(x)$ for all $x$ is also very strong and often not
satisfied. We will be able to say something more precise if we assume risk aversion
$\Rightarrow$ second-order stochastic dominance (SOSD).


\section{Second-order stochastic dominance (SOSD)}

\topic{Definition}
Start with a counterpart to (2): $F$ second-order stochastically dominates $G$ if
\[ \int_0^{x} F(T)\,dT\ \le\ \int_0^{x} G(T)\,dT \qquad \forall\, x\in[0,\bar x]. \tag{3} \]
(2) implies (is stronger than) (3): if $F\le G$ pointwise, the integrals are ordered too.
So we can now compare more distributions.

\topic{A property of (3): a higher expected value}
(3) holds at all $x$, thus at $x=\bar x$. Integrating by parts,
\begin{align*}
  \int_0^{\bar x} 1\cdot F(T)\,dT
   &= \bigl[T\,F(T)\bigr]_0^{\bar x}-\int_0^{\bar x} T\,dF(T)\\
   &= \bar x - 0 - \underbrace{\int_0^{\bar x} T\,dF(T)}_{\text{expected value under } F}.
\end{align*}
So (3) at $\bar x$ reads $\bar x-E_F[x]\le \bar x-E_G[x]$, i.e.\
\[ E_F[x]\ \ge\ E_G[x]: \]
(3) implies a higher (or equal) expected winning under $F$ than under $G$.

\subsection*{Comparing FOSD and SOSD (simple cases)}

\begin{center}
\begin{tikzpicture}[scale=1.25]
  % FOSD: G(t) = t + 0.15 sin(pi t) above F(t) = t - 0.15 sin(pi t); xbar = 4, height 1 = 2.6
  \ecaxes{4.7}{3.1}
  \node[ecnode,right] at (4.7,0) {$x$};
  \draw[dashed] (0,2.6) -- (4,2.6);
  \node[ecnode,left] at (0,2.6) {$1$};
  \draw[dashed] (4,0) node[ecnode,below] {$\bar x$} -- (4,2.6);
  \draw[budget,domain=0:4,samples=60] plot (\x,{2.6*(\x/4+0.15*sin(45*\x))});
  \draw[budget,domain=0:4,samples=60] plot (\x,{2.6*(\x/4-0.15*sin(45*\x))});
  \node[ecnode,anchor=south east] at (1.5,1.55) {$G$};
  \node[ecnode,anchor=north west] at (2.4,1.05) {$F$};
  \node[ecnode] at (2,-0.75) {FOSD $\Rightarrow$ SOSD};
\end{tikzpicture}
\hspace{1.2cm}
\begin{tikzpicture}[scale=1.25]
  % no FOSD but SOSD: G(t) = t + 0.1 sin(pi t), F(t) = t + 0.02 sin(pi t) - 0.09 sin(2 pi t)
  \ecaxes{4.7}{3.1}
  \node[ecnode,right] at (4.7,0) {$x$};
  \draw[dashed] (0,2.6) -- (4,2.6);
  \node[ecnode,left] at (0,2.6) {$1$};
  \draw[dashed] (4,0) node[ecnode,below] {$\bar x$} -- (4,2.6);
  \draw[budget,domain=0:4,samples=80] plot (\x,{2.6*(\x/4+0.1*sin(45*\x))});
  \draw[budget,domain=0:4,samples=80] plot (\x,{2.6*(\x/4+0.02*sin(45*\x)-0.09*sin(90*\x))});
  \node[ecnode,anchor=south east] at (1.3,1.35) {$G$};
  \node[ecnode,anchor=north west] at (1.3,0.45) {$F$};
  \node[ecnode] at (2,-0.75) {no FOSD, but SOSD};
\end{tikzpicture}
\end{center}

\topic{When the CDFs cross: more mass in the tails}

\begin{center}
\begin{tikzpicture}[scale=1.5]
  % F(t) = t (straight line), G(t) = t + 0.12 sin(2 pi t); xbar = 4, height 1 = 2.6
  \ecaxes{4.8}{3.1}
  \node[ecnode,right] at (4.8,0) {$x$};
  \draw[dashed] (0,2.6) -- (4,2.6);
  \node[ecnode,left] at (0,2.6) {$1$};
  \draw[dashed] (4,0) node[ecnode,below] {$\bar x$} -- (4,2.6);
  \draw[budget] (0,0) -- (4,2.6);
  \draw[budget,domain=0:4,samples=80] plot (\x,{2.6*(\x/4+0.12*sin(90*\x))});
  \node[ecnode,anchor=south east] at (1.1,1.12) {$G(x)$};
  \node[ecnode,anchor=north west] at (1.4,0.85) {$F(x)$};
  % x' = 0.8 (t = 0.2): F = 0.52, G = 0.52 + 2.6*0.12*sin(72) = 0.817
  \draw[dashed] (0.8,0) node[ecnode,below] {$x'$} -- (0.8,0.817);
  \draw[dashed] (0,0.817) node[ecnode,left] {$G(x')$} -- (0.8,0.817);
  \draw[dashed] (0,0.52) node[ecnode,left] {$F(x')$} -- (0.8,0.52);
  % x'' = 3.2 (t = 0.8): F = 2.08, G = 2.08 - 0.297 = 1.783
  \draw[dashed] (3.2,0) node[ecnode,below] {$x''$} -- (3.2,2.6);
  \draw[{Stealth}-{Stealth}] (3.1,2.08) -- (3.1,2.6);
  \node[ecnode,left] at (3.1,2.42) {$1-F(x'')$};
  \draw[{Stealth}-{Stealth}] (3.32,1.783) -- (3.32,2.6);
  \node[ecnode,right] at (4.05,2.15) {$1-G(x'')$};
  \draw[thin] (3.36,2.19) -- (4.05,2.15);
\end{tikzpicture}
\end{center}

$\Rightarrow$ $G$ has more mass at the tails ($G(x')>F(x')$ and $1-G(x'')>1-F(x'')$)
$\to$ a heightened probability of very low or very high winnings $\to$ sounds risky.

$\Rightarrow$ Intuition: $F$ is better than $G$ if the agent is risk averse.

\topic{Main result for SOSD}
\[ \int_0^{x} F(T)\,dT\ \le\ \int_0^{x} G(T)\,dT \quad \forall\, x\in[0,\bar x] \tag{3} \]
\[ \Big\Updownarrow \]
\[ \int u(x)\,dF(x)\ \ge\ \int u(x)\,dG(x)
   \quad\text{for every non-decreasing, (weakly) concave } u(x). \tag{4} \]

\topic{Proof of $(3)\Rightarrow(4)$}
Assume $u$ is twice differentiable, so $u'\ge0$ and $u''\le0$. Let $H(x)=F(x)-G(x)$
(so $H(0)=H(\bar x)=0$) and $S(x)=\int_0^{x}H(T)\,dT$, which is $\le0$ for all $x$ by (3).
Integrating by parts twice,
\begin{align*}
  \int u(x)\,dF(x)-\int u(x)\,dG(x)
   &= \int_0^{\bar x} u(x)\,dH(x)
    = \underbrace{\bigl[u(x)H(x)\bigr]_0^{\bar x}}_{=0}-\int_0^{\bar x} u'(x)H(x)\,dx\\
   &= -\Bigl(\bigl[u'(x)S(x)\bigr]_0^{\bar x}-\int_0^{\bar x} u''(x)\,S(x)\,dx\Bigr)\\
   &= -\underbrace{u'(\bar x)}_{\ge0}\,\underbrace{\int_0^{\bar x}H(T)\,dT}_{\le0}
      \;+\;0\;+\;\int_0^{\bar x}\underbrace{u''(x)}_{\le0}\,
      \underbrace{\int_0^{x}H(T)\,dT}_{\le0}\,dx\\
   &\ \ge\ 0.
\end{align*}
(The boundary term at $0$ vanishes because $S(0)=0$.) So (4) holds. $\checkmark$

\topic{Proof of $(4)\Rightarrow(3)$}
A special case of (4) is the (non-decreasing, concave) utility function
\[ u(x)=\begin{cases} x & \text{if } x\le z,\\ z & \text{if } x>z. \end{cases} \]

\begin{center}
\begin{tikzpicture}[scale=0.9]
  \ecaxes{4.4}{2.6}
  \node[ecnode,right] at (4.4,0) {$x$};
  \node[ecnode,above] at (0,2.6) {$u(x)$};
  \draw[budget] (0,0) -- (1.8,1.8) -- (4.2,1.8);
  \draw[dashed] (1.8,0) node[ecnode,below] {$z$} -- (1.8,1.8);
\end{tikzpicture}
\end{center}

Then
\begin{align*}
  \int u(x)\,dF(x) &= \int_0^{z} x\,dF(x)+\int_z^{\bar x} z\,dF(x)\\
   &= \bigl[x\,F(x)\bigr]_0^{z}-\int_0^{z}F(x)\,dx + z\int_z^{\bar x}dF(x)\\
   &= zF(z)-0-\int_0^{z}F(x)\,dx+z\bigl(1-F(z)\bigr)\\
   &= z-\int_0^{z}F(x)\,dx,
\end{align*}
and likewise for $G$. So (4) implies
\[ z-\int_0^{z}F(x)\,dx\ \ge\ z-\int_0^{z}G(x)\,dx
   \quad\Longleftrightarrow\quad \int_0^{z}F(x)\,dx\ \le\ \int_0^{z}G(x)\,dx, \]
and this holds regardless of what $z\in[0,\bar x]$ is $\Rightarrow$ (3). $\checkmark$


\section{General equilibrium: introduction (MWG 15.A--15.C)}

As a micro topic, let us focus primarily on \emph{welfare} questions. There are a lot of
applications in macro and finance.

\begin{flatlist}
  \item We start by using simple models to illustrate important components of the
        analysis and reveal some important conclusions (and pictures).
  \item We bring together consumers, firms and markets, and \emph{endogenize} prices
        and wealth.
  \item We discuss 2 of the 3 models in Chapter 15, omitting the small open economy
        with some exogenous world prices.
\end{flatlist}

\topic{Pure exchange economy}
No production: agents are born with endowments that can be traded.

\runin{Important assumption.} Agents are price takers. Nevertheless, these prices arise
endogenously as part of the equilibrium.

\runin{In equilibrium,} demand $=$ supply in all markets.


\section{An economy with 2 agents and 2 goods}

\begin{flatlist}
  \item Agent $i$ has endowments $\omega_i=(\omega_{1i},\omega_{2i})$, where
        $\omega_{ki}$ is the endowment of good $k$ held by agent $i$ (first index:
        good; second index: agent).
  \item Let $\omega=(\omega_1,\omega_2)$ denote the vector of endowments across agents.
  \item Agent $i$'s bundle is $x_i=(x_{1i},x_{2i})$.
  \item The total endowment of good $k=1,2$ is $\bar\omega_k=\omega_{k1}+\omega_{k2}$.
  \item An allocation $x=(x_1,x_2)$ is \emph{feasible} if
        \[ x_{k1}+x_{k2}\ \le\ \bar\omega_k, \qquad k=1,2. \]
\end{flatlist}

\topic{The Edgeworth box}
We can capture the economy in the Edgeworth box: width $\bar\omega_1$, height
$\bar\omega_2$. Agent 1's origin is the bottom-left corner, agent 2's origin the top-right
corner (agent 2's quantities are measured leftwards and downwards).

\begin{center}
\begin{tikzpicture}[scale=1.0]
  % box 6 x 4; omega = (4, 1.6); budget slope -1.4
  \draw[thick] (0,0) rectangle (6,4);
  \node[ecnode,anchor=north east] at (-0.55,-0.5) {agent 1, $O_1$};
  \node[ecnode,anchor=south west] at (6.55,4.5) {$O_2$, agent 2};
  \draw[dashed] (4,0) -- (4,4);
  \draw[dashed] (0,1.6) -- (6,1.6);
  \draw[budget] (2.286,4) -- (5.143,0);
  \node[ecnode,anchor=south west] at (2.9,2.9) {slope $-\tfrac{p_1}{p_2}$};
  \node[bundle] at (4,1.6) {};
  \node[ecnode,above right] at (4,1.6) {$\omega$};
  % agent 1 measures
  \draw[{Stealth}-{Stealth}] (0,-0.45) -- (4,-0.45);
  \node[ecnode,below] at (2,-0.45) {$\omega_{11}$};
  \draw[{Stealth}-{Stealth}] (-0.45,0) -- (-0.45,1.6);
  \node[ecnode,left] at (-0.45,0.8) {$\omega_{21}$};
  % agent 2 measures
  \draw[{Stealth}-{Stealth}] (6,4.45) -- (4,4.45);
  \node[ecnode,above] at (5,4.45) {$\omega_{12}$};
  \draw[{Stealth}-{Stealth}] (6.45,4) -- (6.45,1.6);
  \node[ecnode,right] at (6.45,2.8) {$\omega_{22}$};
  % box dimensions
  \node[ecnode] at (3,-1.1) {width $\bar\omega_1=\omega_{11}+\omega_{12}$};
  \node[ecnode,rotate=90] at (-1.6,2) {height $\bar\omega_2=\omega_{21}+\omega_{22}$};
\end{tikzpicture}
\end{center}

\begin{flatlist}
  \item Agent $i$'s budget line goes through their endowment, because they do not have
        to trade (or can sell and then buy back their endowment).
  \item Relative prices are the same for both agents, so the two budget lines lie on
        top of each other in the box.
  \item Note that wealth is endogenous: it comes from selling endowments,
        $w_i=p\cdot\omega_i=p_1\omega_{1i}+p_2\omega_{2i}$.
\end{flatlist}


\begin{transcriptnotes}
  \item The export has 7 pages; the lower half of page 7 is blank. No date or title is
        written, so the header date is left empty; the subtitle and the section headings
        for Sections 1--2 and 4 are \textbf{added} (MWG Chapter 6, Section 6.D;
        Chapter 15, Sections 15.A--15.B). The page itself titles Section 3 ``GE ---
        General Equilibrium, Introduction 15.A--15.C''.
  \item Page 1: ``Reminder of FOSD''; the label (1) for the integral condition and (2)
        for $F\le G$ follow Lesson 18. ``not \underline{\ \ } statified'' $\to$
        \emph{often not satisfied} (the blank is filled with ``often'');
        ``what'' is crossed out; ``sat'' $\to$ \emph{say}; ``preceisely'' $\to$
        \emph{precisely}. ``$\Rightarrow$ SOSD'' expanded.
  \item Page 1, (3): written ``$\forall\,T\in[0,\bar x]$''; \textbf{corrected} to
        $\forall\,x\in[0,\bar x]$ ($T$ is the integration variable; page 3 restates (3)
        with $\forall\,x$). ``compare more distributions'' kept; the one-line reason why
        (2) implies (3) is \textbf{added}.
  \item Page 1--2, expected value: ``Exp.\ value of winning under $F$'' with an arrow to
        page 2 ``(3) implies higher exp winning under $F$ than $G$''. The step
        $\bar x-E_F\le\bar x-E_G$ is \textbf{added}; ``higher'' means \emph{weakly}
        higher. The first integrand is written ``$1\,F(T)$'' (the $1$ is the factor
        differentiated by parts); kept as $1\cdot F(T)$.
  \item Page 2, three CDF sketches redrawn to scale; the curves in the first two
        sketches are unlabelled on the page, so the labels $F$, $G$ are \textbf{added}.
        Sketch 1: $G=t+0.15\sin\pi t$, $F=t-0.15\sin\pi t$ ($t=x/\bar x$). Sketch 2
        (curves cross, ``no FOSD, but SOSD''): $G=t+0.1\sin\pi t$,
        $F=t+0.02\sin\pi t-0.09\sin2\pi t$, chosen so that (3) holds with a single
        crossing. Sketch 3: $F$ is the straight line and $G=t+0.12\sin2\pi t$ (a
        mean-preserving spread of $F$). On the page the straight line is drawn past the
        level $1$; it is stopped at $(\bar x,1)$. ``heighten prob'' $\to$
        \emph{heightened probability}; ``winners'' $\to$ \emph{winnings}.
  \item Page 3: ``ntuition'' $\to$ \emph{Intuition}; the page number reads ``3'' with a
        stray ``(4)'' above it. ``For event non-decreasing and (4) weakly concave'' $\to$
        \emph{for every non-decreasing and (weakly) concave}; the ``(4)'' is the
        equation label, moved to the right.
  \item Page 3 proof, $(3)\Rightarrow(4)$: the first line is written
        ``$\int u(x)F(x)$'' $\to$ $\int u(x)\,dF(x)$. In the second integration by parts
        the inner integral is written $\int_0^{0}H(T)\,dT$; \textbf{corrected} to
        $\int_0^{x}H(T)\,dT$. The integral limits $0,\bar x$ and the notation
        $S(x)$ are \textbf{added}; ``$u$ twice differentiable'' and the reasons the
        boundary terms vanish are \textbf{added}. The page marks the signs
        ($u'(\bar x)$: $+$; $\int H$: $-$; $u''$: $-$; $\int H$: $-$) and ``$\ge0$'';
        these are typeset as underbraces.
  \item Page 4, $(4)\Rightarrow(3)$: ``Pf:'' at the top. The utility function is
        $u(x)=\min\{x,z\}$; it is not differentiable at $z$, which is fine because (4)
        is stated for every non-decreasing concave $u$ (the differentiability in the
        previous proof was only used for that direction). The sketch of $u(x)$ is
        redrawn with a flat segment after $z$ (on the page the horizontal axis slopes).
        The parenthesis in ``$z(1-F(z))$'' is closed; ``likewise for $G$'' is
        \textbf{added}.
  \item Page 5: ``primerily'' $\to$ \emph{primarily}; ``a lot of app.'' $\to$
        \emph{applications}; ``anatts and preveal'' $\to$ \emph{analysis and reveal};
        ``+ pics'' $\to$ \emph{(and pictures)}; ``endogenize'' kept. ``dissuc'' $\to$
        \emph{discuss}; ``exogenous world prices'' (``world'' inserted with a caret).
        ``agents are both with endowments'' $\to$ \emph{born with}; ``these prices are
        arise'' $\to$ \emph{these prices arise}.
  \item Page 6 (a fold shows part of the next page, ``(7)'' and ``his'', in the top
        corner): ``eneconomy'' $\to$ \emph{economy}; ``$w_i(w_{1i},w_{2i})$'' $\to$
        $\omega_i=(\omega_{1i},\omega_{2i})$, with the caret note ``Good, agent''
        written out. ``denote the vector of endowments across agents'' and the
        notation $\omega=(\omega_1,\omega_2)$ is \textbf{added}. ``An allocation of
        endowments is feasible'' $\to$ \emph{An allocation $x=(x_1,x_2)$ is feasible}.
  \item Edgeworth box redrawn to scale. On the page the right edge is labelled
        ``length $\bar\omega_1$, height $\bar\omega_2$''; \textbf{corrected} so that the
        width is $\bar\omega_1$ and the height $\bar\omega_2$, with the sums
        $\omega_{k1}+\omega_{k2}$ \textbf{added}. Origins $O_1$, $O_2$ and the dashed
        lines through $\omega$ are \textbf{added}. ``edge worth box'' $\to$
        \emph{Edgeworth box}.
  \item Page 6--7: ``his endowment because he does not have to trade (or can see and
        then but back at his endowment)'' $\to$ \emph{(or can sell and then buy back
        their endowment)}. ``Relative price are \ldots\ budget line are'' $\to$
        \emph{prices are \ldots\ budget lines lie}. ``wealth is endogenous if it comes
        from selling endowments''; the formula $w_i=p\cdot\omega_i$ is \textbf{added}.
\end{transcriptnotes}

\end{document}
